% Segunda lectura — Bloque geométrico: fundamentos y rotaciones (plano, complejos, 3D, eje y exponencial)
\section{Bloque geométrico: fundamentos y rotaciones}

\section{Transformaciones del plano}\label{sec:transf-plano}

Una transformación del plano es una función
\[
  F:\R^2\longrightarrow\R^2.
\]
Estudiarla significa responder cuatro preguntas:
\begin{itemize}
  \item ¿qué puntos permanecen fijos?;
  \item ¿qué magnitudes conserva: distancias, ángulos, áreas o paralelismo?;
  \item ¿cómo se representa algebraicamente?;
  \item ¿qué ocurre al componerla con otra transformación?
\end{itemize}

\subsection{Transformaciones lineales y afines}

Una transformación lineal \(L\) cumple
\[
  L(x+y)=L(x)+L(y),\qquad L(\lambda x)=\lambda L(x).
\]
En particular, necesariamente \(L(0)=0\). Por eso una rotación alrededor del origen puede ser
lineal, pero una rotación con centro distinto del origen no puede serlo en esas coordenadas.

Una transformación afín agrega una traslación:
\[
  F(x)=Ax+t.
\]
La matriz \(A\) determina la deformación lineal y el vector \(t\) desplaza el resultado.

\section{Rotaciones en el plano}\label{sec:rot-plano}

\subsection{Matriz de rotación}

La rotación antihoraria de ángulo \(\theta\) alrededor del origen es
\[
  R(\theta)=
  \begin{pmatrix}
    \cos\theta&-\sin\theta\\
    \sin\theta& \cos\theta
  \end{pmatrix}.
\]
Aplicada a \(x=(x_1,x_2)^T\):
\[
  R(\theta)x=
  \begin{pmatrix}
    x_1\cos\theta-x_2\sin\theta\\
    x_1\sin\theta+x_2\cos\theta
  \end{pmatrix}.
\]

\begin{center}
\begin{tikzpicture}[scale=1.0]
  \draw[axisline,-{Latex[length=2mm]}] (-.3,0)--(5.2,0) node[right] {$x$};
  \draw[axisline,-{Latex[length=2mm]}] (0,-.3)--(0,4.4) node[above] {$y$};
  \coordinate (p) at (4,1);
  \coordinate (q) at (2.2,3.5);
  \draw[auxline] (0,0)--(p);
  \draw[auxline] (0,0)--(q);
  \draw[mainline,-{Latex[length=2.6mm]}] (p)
    arc[start angle=14,end angle=58,radius=4.12];
  \fill[teal] (p) circle (2.5pt) node[below right] {$x$};
  \fill[rust] (q) circle (2.5pt) node[above right] {$R(\theta)x$};
  \node[text=teal] at (2.1,.75) {$\theta$};
\end{tikzpicture}
\end{center}

\begin{resultado}[Propiedades]
\[
R(\theta)^TR(\theta)=I,\qquad
\det R(\theta)=1,\qquad
R(\alpha)R(\beta)=R(\alpha+\beta).
\]
La ortogonalidad implica que conserva producto interno, longitudes y ángulos. El determinante
positivo indica que también conserva la orientación.
\end{resultado}

\paragraph{Ejemplo.}
Para \(\theta=90^\circ\):
\[
R\left(\frac{\pi}{2}\right)=
\begin{pmatrix}0&-1\\1&0\end{pmatrix},
\qquad
R\left(\frac{\pi}{2}\right)(3,1)^T=(-1,3)^T.
\]

\subsection{Rotación alrededor de un centro cualquiera}

Para rotar alrededor de \(c\), primero se traslada el centro al origen, luego se rota y
finalmente se deshace la traslación:
\[
  \boxed{Q_{c,\theta}(x)=c+R(\theta)(x-c)}.
\]
Equivalentemente,
\[
  Q_{c,\theta}(x)=R(\theta)x+\bigl(I-R(\theta)\bigr)c.
\]
Es una transformación afín. Si \(c\neq0\), el término independiente suele ser distinto de
cero y la transformación no es lineal.

\begin{center}
\begin{tikzpicture}[scale=.95]
  \coordinate (c) at (3,2.2);
  \coordinate (p) at (6.2,1.5);
  \coordinate (q) at (5.4,4.3);
  \draw[auxline] (c)--(p);
  \draw[auxline] (c)--(q);
  \draw[mainline,-{Latex[length=2.6mm]}] (p)
    arc[start angle=-12,end angle=42,radius=3.28];
  \fill[ink] (c) circle (2.7pt) node[below left] {$c$};
  \fill[teal] (p) circle (2.5pt) node[below right] {$x$};
  \fill[rust] (q) circle (2.5pt) node[above right] {$Q_{c,\theta}(x)$};
  \draw[muted,dotted] (0,0)--(7.6,0);
  \draw[muted,dotted] (0,0)--(0,5);
  \fill[muted] (0,0) circle (1.8pt) node[below left] {$0$};
\end{tikzpicture}
\end{center}

\section{Rotaciones y números complejos}\label{sec:rot-complejos}

Identificando \((x,y)\in\R^2\) con \(z=x+iy\in\mathbb C\), la fórmula de Euler
\[
  e^{i\theta}=\cos\theta+i\sin\theta
\]
permite escribir una rotación como
\[
  z'=e^{i\theta}z.
\]
Si se multiplica por un complejo general \(\rho e^{i\theta}\), el módulo \(\rho\) produce
una escala uniforme y el argumento \(\theta\) produce una rotación:
\[
  z'=\rho e^{i\theta}z.
\]

\begin{idea}[Relación entre la forma compleja y la matricial]
\[
(a+ib)(x+iy)=(ax-by)+i(bx+ay)
\]
equivale a
\[
\begin{pmatrix}x'\\y'\end{pmatrix}
=
\begin{pmatrix}a&-b\\b&a\end{pmatrix}
\begin{pmatrix}x\\y\end{pmatrix}.
\]
La multiplicación compleja es una familia especial de transformaciones lineales reales:
solo usa dos parámetros, frente a los cuatro de una matriz real \(2\times2\) general.
\end{idea}

\section{Rotaciones en tres dimensiones}\label{sec:rot-3d}

\subsection{Rotaciones elementales}

Las rotaciones alrededor de los ejes coordenados se representan por
\[
R_x(\alpha)=
\begin{pmatrix}
1&0&0\\0&\cos\alpha&-\sin\alpha\\0&\sin\alpha&\cos\alpha
\end{pmatrix},
\]
\[
R_y(\beta)=
\begin{pmatrix}
\cos\beta&0&\sin\beta\\0&1&0\\-\sin\beta&0&\cos\beta
\end{pmatrix},
\qquad
R_z(\gamma)=
\begin{pmatrix}
\cos\gamma&-\sin\gamma&0\\
\sin\gamma& \cos\gamma&0\\
0&0&1
\end{pmatrix}.
\]

\subsection{El orden importa}\label{sec:orden-importa}

En general,
\[
  R_x(\alpha)R_y(\beta)R_z(\gamma)
  \neq
  R_z(\gamma)R_y(\beta)R_x(\alpha).
\]
La matriz situada más a la derecha actúa primero. Esta falta de conmutatividad fue uno de los
experimentos principales de la primera clase.

\paragraph{Ejemplo sencillo.}
Con giros de \(90^\circ\), aplicar primero \(R_x\) y luego \(R_y\) al vector
\(e_z=(0,0,1)^T\) produce
\[
R_y\left(\frac{\pi}{2}\right)
R_x\left(\frac{\pi}{2}\right)e_z
=R_y\left(\frac{\pi}{2}\right)(-e_y)
=-e_y.
\]
Al invertir el orden:
\[
R_x\left(\frac{\pi}{2}\right)
R_y\left(\frac{\pi}{2}\right)e_z
=R_x\left(\frac{\pi}{2}\right)e_x
=e_x.
\]
Como \(-e_y\neq e_x\), la diferencia no es un error numérico: es una propiedad geométrica
de las rotaciones espaciales.

\subsection{Rotaciones pequeñas}

Para un ángulo pequeño \(\varepsilon\),
\[
  R_x(\varepsilon)\approx I+\varepsilon A_x,
  \qquad
  R_y(\varepsilon)\approx I+\varepsilon A_y.
\]
Entonces
\[
(I+\varepsilon A_x)(I+\varepsilon A_y)
=I+\varepsilon(A_x+A_y)+\varepsilon^2A_xA_y,
\]
mientras que el orden inverso cambia solo el término de segundo orden:
\[
(I+\varepsilon A_y)(I+\varepsilon A_x)
=I+\varepsilon(A_x+A_y)+\varepsilon^2A_yA_x.
\]
Por eso las rotaciones muy pequeñas parecen conmutar a primer orden, aunque no conmuten
exactamente.

\section{Eje de rotación, autovectores y exponencial}\label{sec:eje-exp}

\subsection{El eje como dirección fija}

Toda rotación propia de \(\R^3\) posee un eje. Si \(v\) apunta en la dirección de ese eje,
\[
  Rv=v.
\]
Por lo tanto, \(v\) es un autovector de \(R\) con autovalor \(1\). Equivalentemente,
\[
  (R-I)v=0.
\]

\subsection{Exponencial de matrices}

La exponencial matricial se define por la serie
\[
  e^A=I+A+\frac{A^2}{2!}+\frac{A^3}{3!}+\cdots.
\]
Si una rotación se expresa como \(R=e^A\), entonces
\[
  R^t=e^{tA},\qquad 0\le t\le1,
\]
describe una trayectoria continua desde la identidad hasta \(R\). Del mismo modo,
\[
  e^{A/n}
\]
es un paso pequeño que, aplicado \(n\) veces, produce \(e^A\).

\begin{resultado}[Generador antisimétrico]
Para una rotación tridimensional, el logaritmo matricial principal produce, bajo condiciones
adecuadas, una matriz antisimétrica
\[
A=
\begin{pmatrix}
0&-\omega_z&\omega_y\\
\omega_z&0&-\omega_x\\
-\omega_y&\omega_x&0
\end{pmatrix}.
\]
El vector \(\omega=(\omega_x,\omega_y,\omega_z)\) indica el eje y su norma representa el
ángulo de rotación.
\end{resultado}

\subsection{Experimento de la primera clase}

Se construyeron en Excel o LibreOffice las matrices \(R_x,R_y,R_z\) y se compararon distintos
órdenes con la función \texttt{MMULT}. Después se aplicó muchas veces una composición de
giros pequeños e iguales alrededor de los tres ejes. El resultado aproxima una rotación
alrededor de la dirección \((1,1,1)\), capaz de permutar cíclicamente los ejes:
\[
  e_x\longmapsto e_y,\qquad
  e_y\longmapsto e_z,\qquad
  e_z\longmapsto e_x.
\]
El ángulo de esa permutación es \(120^\circ\), no \(90^\circ\).

Para cuantificar el experimento, puede tomarse un paso
\[
Q_\varepsilon
=R_z(\varepsilon)R_y(\varepsilon)R_x(\varepsilon),
\qquad \varepsilon=0.01.
\]
A primer orden,
\[
Q_\varepsilon
\approx I+\varepsilon(A_x+A_y+A_z).
\]
La suma de generadores corresponde al eje \((1,1,1)\), cuya norma es \(\sqrt3\). Después de
\(n\) pasos, el ángulo acumulado es aproximadamente
\[
  \Theta\approx n\varepsilon\sqrt3.
\]
Con \(n=121\):
\[
  \Theta\approx121(0.01)\sqrt3
  \approx2.096\ {\rm rad}
  \approx120.1^\circ.
\]
La rotación exacta de \(120^\circ\) alrededor de \((1,1,1)\) que realiza la permutación
anterior tiene matriz
\[
P=
\begin{pmatrix}
0&0&1\\
1&0&0\\
0&1&0
\end{pmatrix}.
\]
Las columnas muestran directamente que \(Pe_x=e_y\), \(Pe_y=e_z\) y \(Pe_z=e_x\).

\begin{idea}[Práctica recomendada]
\begin{enumerate}
  \item Construir \(R_x,R_y,R_z\) en una hoja de cálculo.
  \item Comparar dos órdenes con ángulos grandes.
  \item Repetir con \(\varepsilon=0.01\) radianes.
  \item Calcular en Python \(\log(R)\) con \texttt{scipy.linalg.logm}.
  \item Interpretar la matriz antisimétrica obtenida.
\end{enumerate}
\end{idea}

